% !TEX root =  ../Dissertation.tex

\chapter{Experiments}

This chapter presents a comprehensive experimental study on SER under clean and noisy conditions. We evaluate multiple architectures and noise augmentation strategies on two datasets, adhering to the unified notation and preprocessing defined in Chapter 3. Unless otherwise specified, audio is sampled at \(f_s\), the time-domain signal is \(x[n]\), noise is \(v[n]\), and on-the-fly mixtures during training follow
\(x_{\mathrm{mix}}[n] = x_s[n] + g\, x_n[n]\) with gain \(g\) computed from target SNR as in Chapter 3.

\section{Common Settings}

\subsection{Datasets and splits}
We use RAVDESS (speech actors 01–24) and IEMOCAP. Data splits follow the established protocol for each dataset. Validation and clean test sets are never augmented; only the training split is subjected to on-the-fly noise mixing when enabled. The batch size of MLP, shallow CNN, and Resnet18 is 128, and 64 for AST. All models are trained using Adam \cite{Kingma2014Adam} optimizer with a initial learning rate of 0.001. The training is terminated when the validation loss does not improve for 10 epochs.

\paragraph{Class schemes and dataset sizes.}
For categorical SER we adopt fixed emotion inventories per corpus.
\begin{itemize}
  \item \textbf{RAVDESS (8 classes).} Neutral, Calm, Happy, Sad, Angry, Fearful, Disgust, Surprised. We use the speech subset from actors 01–24 and retain all eight categories. Splits are speaker-disjoint following the common protocol. Dataset sizes used in this work (utterances): \textit{Train = 1152,  Test = 288}.
  \item \textbf{IEMOCAP (4 classes).} We adopt the standard 4-class mapping with Happy and Excited merged into Happy: Angry, Happy, Sad, Neutral. Sessions are used for speaker-disjoint splits. Dataset sizes (utterances): \textit{Train = 10560, Test = 2896}.
\end{itemize}
These label spaces remain fixed across Experiments A–D. All reported metrics are computed on the same class definitions; no relabeling is performed between experiments.

\subsection{Models}
We compare a convolutional baseline (ResNet-18 on Log-Mel spectrograms) and a transformer-based model (AST Base on filterbank features), alongside lighter baselines where applicable (MLP and shallow CNN) to contextualize complexity–performance trade-offs.

\paragraph{Implementation summary.}
Models are instantiated via a configuration-driven factory in \texttt{models.py}: \texttt{create\_model(config)} reads \texttt{config['training']}

\texttt{['model\_name']} and \texttt{config['models'][...]['type']} to build the requested architecture (MLP, ShallowCNN, ResNet18, AST). Core hyperparameters—including \textit{num\_classes}, dropout rates, hidden sizes, and feature types—are taken from the configuration for reproducibility and easy switching between architectures.

\begin{itemize}
  \item \textbf{MLP (MFCC-based).} Input features are MFCC statistics (per-utterance mean and standard deviation concatenated; dimension \(2\times\textit{input\_size}\)). The backbone is a sequence of \texttt{Linear→ReLU→Dropout} blocks; a final linear classifier produces emotion logits. This baseline targets low-dimensional hand-crafted features with small parameter count and stable training.
  \item \textbf{Shallow CNN (Log-Mel-based).} The convolutional trunk has three blocks with channels \(32\to64\to128\). Each block applies \texttt{Conv2d + BatchNorm + ReLU + MaxPool}; an adaptive global average pooling aggregates time–frequency features. A dropout layer and a single \texttt{Linear} head complete classification. The design is lightweight yet effective for 2D time–frequency modeling.
  \item \textbf{ResNet-18 (Log-Mel-based).} The network begins with a \(7\times7\) convolution and max-pooling, followed by four residual stages (channels \(64/128/256/512\), increasing strides). Each \textit{BasicBlock} uses two \(3\times3\) convolutions with identity/projection shortcuts. After adaptive average pooling, a fully connected layer outputs logits. Kaiming initialization and BatchNorm parameter initialization are used; \texttt{num\_classes} is validated at construction.
  \item \textbf{AST (Audio Spectrogram Transformer).} Implemented in \texttt{ast\_model.py} using a DeiT/ViT variant (via \texttt{timm}). \textit{PatchEmbed} is adapted for single-channel Log-Mel; 3-channel weights are merged to 1-channel, and patch grid plus positional embeddings are recomputed (bilinear interpolation or cropping) to match audio shapes. ImageNet pretraining is supported (optionally AudioSet). The head is \texttt{LayerNorm + Linear}. A wrapper (e.g., \texttt{RAVDESSASTModel}) derives input time steps from sampling rate, duration, and \(n\_\mathrm{mels}\), instantiating AST accordingly. Forward path: reshape \((B,T,F)\to(B,1,F,T)\), apply patch embedding, add \texttt{[cls]/[dist]} tokens and positional encodings, pass transformer blocks and normalization, then pool \texttt{cls}/\texttt{dist} representations and use an MLP head to obtain logits.
\end{itemize}



\subsection{Training and evaluation}
Training uses mini-batch stochastic optimization with standard regularization. Metrics on the validation set include accuracy, precision, recall, and macro-F1. Confusion matrices are reported to analyze class-wise separability. Robustness is further assessed on noisy test sets created by offline mixing of \(x[n]\) with held-out noise at specified SNRs \(\in \{0,5,10,20\}\,\mathrm{dB}\). All experiments fix a base random seed 42 unless otherwise stated.

\subsection{Noise augmentation policy}
When enabled, augmentation follows Chapter 3: with probability \(p_{\mathrm{aug}}\), we sample a noise type (white or ESC-50 subsets), sample SNR from \(\{0,5,10,20\}\,\mathrm{dB}\), time-align noise to the utterance, compute gain \(g\) via RMS matching to satisfy the target SNR, form \(x_{\mathrm{mix}}[n]\), apply peak limiting at \(-1\,\mathrm{dBFS}\), and optionally mean–variance normalize prior to feature extraction.

\section{Experiment A: Clean Baselines}
\textbf{Objective.} Establish upper bounds under matched clean conditions to isolate modeling capacity from noise robustness. These baselines indicate how well each architecture fits clean speech and provide a reference point for subsequent robustness interventions.

\textbf{Setup.} All training, validation, and test audio remain unaltered. Feature pipelines follow Chapter 3 (MFCC and log-mel) as required by each model. Optimization hyperparameters are fixed across runs for comparability. No augmentation, denoising, or calibration is applied. Training–validation trajectories are examined to gauge overfitting and learning dynamics.

\textbf{Evaluation.} We report accuracy, precision, recall, and macro-F1, with confusion matrices to visualize class confusability. Macro-F1 is emphasized due to class imbalance, ensuring that performance is not dominated by frequent emotions. Where informative, we comment on calibration and thresholding effects.

\textbf{Results and analysis.} Results of Experiment A on RAVDESS and IEMOCAP are shown in Table \ref{tab:ravdess_results} and Table \ref{tab:iemocap_results} respectively. Figure \ref{fig:expA_barplots} provides a more intuitive view of the results. On RAVDESS, AST Base achieves the best performance (Acc 0.718, Macro-F1 0.714), closely followed by ResNet-18 (Acc 0.687). MLP offers a moderate reference (Acc 0.451), while the shallow CNN underfits substantially (Acc 0.208), suggesting inadequate receptive-field depth for clean spectral patterns. On IEMOCAP, AST again leads in accuracy (0.615) but exhibits a precision–recall asymmetry (Precision 0.682 vs. Recall 0.413), indicative of conservative predictions that miss minority or ambiguous classes. ResNet-18 is competitive yet shows greater variability across sessions, and MLP/CNN display sensitivity to the dataset’s heterogeneity and longer utterances. Confusions concentrate among emotions with overlapping prosody (e.g., neutral–sad, angry–excited), consistent with prior findings. Training–validation comparisons suggest mild overfitting for deeper models on RAVDESS and larger variance on IEMOCAP, pointing to limited per-class data. These ceilings highlight headroom primarily in recall on IEMOCAP and in narrowing the AST–ResNet gap on RAVDESS.





\begin{table}[htbp]
    \centering
    \caption{ Results of Experiment A on RAVDESS}
    \label{tab:ravdess_results}
    \begin{tabular}{lcccc}
        \hline
        Model & Accuracy & Precision & Recall & Macro-F1 \\
        \hline
        MLP & 0.451 & 0.412 & 0.431 & 0.402 \\
        Shallow CNN & 0.208 & 0.141 & 0.192 & 0.127 \\
        ResNet-18 & 0.687 & 0.694 & 0.697 & 0.682 \\
        AST Base &  0.718 &  0.735 & 0.719 & 0.714 \\
        \hline
    \end{tabular}
\end{table}

\begin{table}[htbp]
    \centering
    \caption{ Results of Experiment A on IEMOCAP}
    \label{tab:iemocap_results}
    \begin{tabular}{lcccc}
        \hline
        Model & Accuracy & Precision & Recall & Macro-F1 \\
        \hline
        MLP & 0.583 & 0.498 & 0.342 & 0.340 \\
        Shallow CNN & 0.583 & 0.546 & 0.449 & 0.415 \\
        ResNet-18 & 0.570 & 0.388 & 0.475 & 0.406 \\
        AST Base & 0.615 & 0.682 & 0.413 &  0.448\\

        \hline
    \end{tabular}
\end{table}


\begin{figure}[htbp]
    \centering
    \begin{tikzpicture}
    \pgfplotsset{compat=1.18}

    \begin{groupplot}[
      group style={group size=2 by 1, horizontal sep=4.5em},
      width=0.65\textwidth, height=0.60\textwidth,
      ymin=0, ymax=0.85,
      ymajorgrids, grid=both,
      grid style={line width=.1pt, draw=gray!30},
      major grid style={line width=.2pt, draw=gray!55},
      xtick=data,
      symbolic x coords={MLP, Shallow CNN, ResNet-18, AST Base},
      xticklabel style={font=\tiny, rotate=20, anchor=east},
      ylabel={Score},
      enlarge x limits=0.22,
      ytick distance=0.1,
      ticklabel style={font=\small},
      label style={font=\small}
    ]

    % ---------------- RAVDESS ----------------
    \nextgroupplot[
      title={RAVDESS}, title style={font=\small},
      ybar,
      bar width=9pt,
      legend style={
        at={(1.5,1.25)}, anchor=south,
        nodes={scale=0.9, transform shape},
        legend columns=4, /tikz/every even column/.style={column sep=0.8em},
        row sep=2pt
      }
    ]
    % Accuracy
    \addplot+[draw=teal!60!black, fill=teal!80!white,area legend,
    %   every node near coord/.append style={font=\scriptsize, /pgf/number format/fixed, /pgf/number format/precision=2, yshift=1.5pt}
    ] coordinates
    {(MLP,0.451) (Shallow CNN,0.208) (ResNet-18,0.687) (AST Base,0.718)};
    \addlegendentry{Accuracy}
    % Precision
    \addplot+[draw=magenta!60!black, fill=magenta!80!white,area legend,
      every node near coord/.append style={font=\scriptsize, /pgf/number format/fixed, /pgf/number format/precision=2, yshift=1.5pt}
    ] coordinates
    {(MLP,0.412) (Shallow CNN,0.141) (ResNet-18,0.694) (AST Base,0.735)};
    \addlegendentry{Precision}
    % Recall
    \addplot+[draw=orange!70!black, fill=orange!90!white,area legend,
      every node near coord/.append style={font=\scriptsize, /pgf/number format/fixed, /pgf/number format/precision=2, yshift=1.5pt}
    ] coordinates
    {(MLP,0.431) (Shallow CNN,0.192) (ResNet-18,0.697) (AST Base,0.719)};
    \addlegendentry{Recall}
    % Macro-F1
    \addplot+[ draw=violet!70!black, fill=violet!85!white,area legend,
      every node near coord/.append style={font=\scriptsize, /pgf/number format/fixed, /pgf/number format/precision=2, yshift=1.5pt}
    ] coordinates
    {(MLP,0.402) (Shallow CNN,0.127) (ResNet-18,0.682) (AST Base,0.714)};
    \addlegendentry{Macro-F1}

    % ---------------- IEMOCAP ----------------
    \nextgroupplot[title={IEMOCAP}, title style={font=\small},      ybar,
    bar width=9pt,
    ]
    % Accuracy
    \addplot+[ draw=teal!60!black, fill=teal!80!white, area legend,
      every node near coord/.append style={font=\scriptsize, /pgf/number format/fixed, /pgf/number format/precision=2, yshift=1.5pt}
    ] coordinates
    {(MLP,0.583) (Shallow CNN,0.583) (ResNet-18,0.570) (AST Base,0.615)};
    % Precision
    \addplot+[ draw=magenta!60!black, fill=magenta!80!white,area legend,
      every node near coord/.append style={font=\scriptsize, /pgf/number format/fixed, /pgf/number format/precision=2, yshift=1.5pt}
    ] coordinates
    {(MLP,0.498) (Shallow CNN,0.546) (ResNet-18,0.388) (AST Base,0.682)};
    % Recall
    \addplot+[ draw=orange!70!black, fill=orange!90!white,area legend,
      every node near coord/.append style={font=\scriptsize, /pgf/number format/fixed, /pgf/number format/precision=2, yshift=1.5pt}
    ] coordinates
    {(MLP,0.342) (Shallow CNN,0.449) (ResNet-18,0.475) (AST Base,0.413)};
    % Macro-F1
    \addplot+[ draw=violet!70!black, fill=violet!85!white,area legend,
      every node near coord/.append style={font=\scriptsize, /pgf/number format/fixed, /pgf/number format/precision=2, yshift=1.5pt}
    ] coordinates
    {(MLP,0.340) (Shallow CNN,0.415) (ResNet-18,0.406) (AST Base,0.448)};

    \end{groupplot}
    \end{tikzpicture}
    \caption{Experiment A: Grouped bar charts of four models across four metrics on RAVDESS and IEMOCAP.}
    \label{fig:expA_barplots}
    \end{figure}


\section{Experiment B: White-Noise Training at 20 dB}


\textbf{Objective.} Test whether injecting mild additive white noise during training (\(\mathrm{SNR}=20\,\mathrm{dB}\)) enhances robustness while preserving accuracy on clean evaluation. We also examine how this augmentation shifts the precision–recall balance and class confusability relative to Experiment A.

\textbf{Setup.} During training, white noise at \(\mathrm{SNR}=20\,\mathrm{dB}\) is mixed with probability \(p_{\mathrm{aug}}\) using RMS-based gain control (Chapter 3). Validation and the primary test set remain clean to isolate changes in clean-set generalization. Hyperparameters and feature configurations are held fixed across experiments for fair comparison.

\textbf{Evaluation.} We report accuracy, precision, recall, and macro-F1, emphasizing macro-F1 to account for class imbalance. Results are compared directly against Experiment A to quantify any trade-offs introduced by augmentation.

\textbf{Results and analysis.} Table \ref{tab:white_noise_results} shows the results of Experiment B on RAVDESS and IEMOCAP. On RAVDESS, clean-set performance is essentially preserved for AST Base (Acc 0.718, Macro-F1 0.713 vs. 0.718/0.714 in Exp A), indicating that 20 dB noise does not harm its clean discriminability. ResNet-18 exhibits a small decline (Acc 0.677, Macro-F1 0.674 vs. 0.687/0.682), consistent with a mild regularization effect that slightly shifts decision boundaries without substantial loss. On IEMOCAP, effects are model-dependent. ResNet-18 shows higher precision (0.491 vs. 0.388) but lower recall (0.378 vs. 0.475) and macro-F1 (0.315 vs. 0.406), suggesting a more conservative operating regime that reduces false positives at the expense of coverage, potentially reflecting session and speaker variability. In contrast, AST Base improves recall (0.471 vs. 0.413) and macro-F1 (0.472 vs. 0.448) while experiencing lower accuracy (0.583 vs. 0.615) and precision (0.565 vs. 0.682). This pattern indicates better balance across classes—even if top-1 correctness declines—consistent with augmentation helping the transformer capture noise-invariant cues on a heterogeneous corpus. Overall, training with 20 dB white noise acts as a gentle regularizer: it preserves clean-set performance on the simpler RAVDESS distribution and yields dataset-dependent precision–recall shifts on IEMOCAP, with net gains in macro-F1 for AST and modest degradations for ResNet-18.

\begin{table}[htbp]
    \centering
    \caption{ Results of Experiment B on two datasets }
    \label{tab:white_noise_results}
    \setlength{\tabcolsep}{5pt}
    \resizebox{\textwidth}{!}{%
    \begin{tabular}{lcccccccc}
        \hline
         & \multicolumn{4}{c}{RAVDESS} & \multicolumn{4}{c}{IEMOCAP} \\
        Model & Acc. & Prec. & Recall & Macro-F1 & Acc. & Prec. & Recall & Macro-F1 \\
        \hline
        ResNet-18 & 0.677 & 0.674 & 0.680 & 0.674 &  0.551 &  0.491 & 0.378 & 0.315 \\
        AST Base & 0.718 & 0.717 & 0.716 & 0.713 & 0.583 & 0.565 & 0.4713 & 0.472 \\
        \hline
    \end{tabular}%
    }
\end{table}



\section{Experiment C: ESC-50 Multi-Category, Multi-SNR}
\textbf{Objective.} Evaluate whether multi-category, multi-SNR environmental noise augmentation (ESC-50; natural soundscapes and human non-speech, \(\mathrm{SNR}\in\{0,5,10,20\}\,\mathrm{dB}\)) builds noise-invariant representations while preserving or improving clean-set performance. We compare architectures and datasets to quantify precision–recall shifts under broader nuisance variability.

\textbf{Setup.} During training, a noise category and SNR are sampled uniformly, and mixed via RMS-controlled gain with peak limiting at \(-1\,\mathrm{dBFS}\) (Chapter 3). Validation and the primary test sets remain clean; noisy test sets are additionally constructed and stratified by category and SNR. Feature pipelines and hyperparameters are kept fixed for fair comparison across experiments.

\textbf{Evaluation.} We report accuracy, precision, recall, and macro-F1, emphasizing macro-F1 to mitigate class-imbalance effects. Results are contrasted with Experiments A/B to isolate the impact of realistic, heterogeneous augmentation from white-noise regularization.

\textbf{Results and analysis.} Table \ref{tab:esc_noise_results} shows the results of Experiment C on RAVDESS and IEMOCAP. On RAVDESS, AST Base improves over the clean baseline (Acc 0.739 vs. 0.718; Macro-F1 0.734 vs. 0.714), indicating that diverse ESC-50 perturbations enhance generalization without harming clean discriminability. In contrast, ResNet-18 degrades notably (Acc 0.548; Macro-F1 0.527 vs. 0.687/0.682 in Exp A), suggesting sensitivity to the broader noise distribution and possible misalignment between convolutional inductive biases and heterogeneous background cues. On IEMOCAP, AST Base again benefits: recall rises (0.488 vs. 0.413) and macro-F1 increases (0.515 vs. 0.448), with a slight accuracy gain (0.625 vs. 0.615) and modest precision reduction (0.645 vs. 0.682). This pattern reflects improved coverage of minority or ambiguous classes, consistent with more balanced decision boundaries under augmentation. By contrast, ResNet-18 exhibits a pronounced precision–recall trade-off (precision 0.519↑, recall 0.327↓; macro-F1 0.323↓ from 0.406), indicating a conservative operating regime that reduces false positives at the expense of sensitivity. Overall, multi-SNR, multi-category augmentation appears especially effective for the transformer (AST) architecture, yielding higher ceilings on the simpler RAVDESS distribution and recall-oriented gains on the heterogeneous IEMOCAP corpus, while the CNN (ResNet-18) shows vulnerability to over-regularization or distributional mismatch.


\begin{table}[htbp]
    \centering
    \caption{ Results of Experiment C on two datasets }
    \label{tab:esc_noise_results}
    \setlength{\tabcolsep}{5pt}
    \resizebox{\textwidth}{!}{%
    \begin{tabular}{lcccccccc}
        \hline
         & \multicolumn{4}{c}{RAVDESS} & \multicolumn{4}{c}{IEMOCAP} \\
        Model & Acc. & Prec. & Recall & Macro-F1 & Acc. & Prec. & Recall & Macro-F1 \\
        \hline
        ResNet-18 & 0.548 & 0.541 & 0.546 & 0.527 & 0.554 & 0.519 & 0.327 & 0.323 \\
        AST Base & 0.739 & 0.746 & 0.745 & 0.734 & 0.625 & 0.645 & 0.488 & 0.515 \\
        \hline
    \end{tabular}%
    }
\end{table}

\section{Experiment D: SNR Ablation (Natural Soundscapes)}
\textbf{Objective.} Evaluate whether multi-category, multi-SNR environmental noise augmentation (ESC-50; natural soundscapes and human non-speech, \(\mathrm{SNR}\in\{0,5,10,20\}\,\mathrm{dB}\)) builds noise-invariant representations while preserving or improving clean-set discriminability. The experiment compares architectures and datasets to quantify shifts in accuracy, precision–recall balance, and macro-F1 under heterogeneous nuisance variability, and to assess whether the learned decision boundaries become more class-balanced than those obtained in Experiments A/B.

\textbf{Setup.} During training, a noise category and an SNR are sampled uniformly and mixed through RMS-controlled gain with peak limiting at \(-1\,\mathrm{dBFS}\) (Chapter 3). Validation and the primary test sets remain clean to isolate changes in clean-set generalization; additional noisy test sets are constructed and stratified by category and SNR for robustness evaluation. Feature pipelines and optimization hyperparameters are held fixed across experiments to ensure comparability.

\textbf{Evaluation.} We report accuracy, precision, recall, and macro-F1. Macro-F1 is emphasized to mitigate class-imbalance artifacts and to reflect model behavior on minority and ambiguous emotions. Qualitative analyses use confusion matrices and heatmaps (model × noise × SNR) to visualize boundary shifts and category-specific confusability. Aggregate results are compared with Experiments A (clean baselines) and B (white-noise at 20 dB) to contextualize the impact of realistic environmental noise.

\textbf{Results and analysis.} Table \ref{tab:snr_ablation_ravdess} and Table \ref{tab:snr_ablation_iemocap} show the results of Experiment D on RAVDESS and IEMOCAP respectively. Line charts \ref{fig:snr_ablation_macroF1} provide a more intuitive view of the results. On RAVDESS, AST Base surpasses the clean baseline, achieving Acc 0.739 and Macro-F1 0.734 (vs. 0.718/0.714 in Experiment A). Precision and recall remain well balanced (0.746/0.745), indicating that broad ESC-50 perturbations help AST maintain clean discriminability while enhancing coverage across classes. In contrast, ResNet-18 shows a marked decrease to Acc 0.548 and Macro-F1 0.527 (vs. 0.687/0.682), with precision and recall (0.541/0.546) remaining roughly symmetric but at a lower operating point. This gap suggests that the transformer benefits from diverse temporal–spectral cues introduced by realistic noise, whereas the convolutional model is more sensitive to the heterogeneity and may over-regularize away discriminative fine structure on the simpler RAVDESS distribution.

On IEMOCAP, AST Base again improves relative to Experiment A: Acc increases to 0.625 (from 0.615), recall rises to 0.488 (from 0.413), and Macro-F1 reaches 0.515 (from 0.448), while precision decreases moderately to 0.645 (from 0.682). The pattern reflects a more balanced decision surface that captures minority or ambiguous emotions without substantial loss of overall correctness. ResNet-18, however, shifts to a conservative regime: precision increases to 0.519 (from 0.388), but recall drops to 0.327 (from 0.475), bringing Macro-F1 down to 0.323 (from 0.406). This precision–recall asymmetry indicates tighter boundaries that reduce false positives at the cost of sensitivity, a behavior consistent with session and speaker variability interacting with heterogeneous noise.


Relative to white-noise training (Experiment B), multi-category augmentation yields clearer benefits for AST. On RAVDESS, AST’s Macro-F1 increases further (0.734 vs. 0.713), and on IEMOCAP, Macro-F1 improves as well (0.515 vs. 0.472), with gains primarily driven by higher recall. For ResNet-18, the multi-category setting underperforms white-noise on RAVDESS (Macro-F1 0.527 vs. 0.674), while on IEMOCAP its Macro-F1 is similar in magnitude (0.323 vs. 0.315) but exhibits a stronger precision bias. Taken together, these results show that multi-SNR, multi-category augmentation is especially effective for the transformer architecture, which maintains or raises clean-set performance on RAVDESS and improves recall and macro-F1 on the more heterogeneous IEMOCAP corpus. The CNN, by contrast, displays model-dependent trade-offs and a greater susceptibility to the diversity of environmental noise, as reflected by reduced macro-F1 and altered calibration on clean evaluation.


\begin{table}[htbp]
    \centering
    \caption{Experiment D (RAVDESS): SNR ablation with natural soundscapes for ResNet-18 and AST. Metrics reported are accuracy (Acc.) and macro-F1 per SNR.}
    \label{tab:snr_ablation_ravdess}
    \begin{tabular}{lcccc}
        \hline
        SNR (dB) & ResNet-18 Acc. & ResNet-18 Macro-F1 & AST Acc. & AST Macro-F1 \\
        \hline
        0  & 0.625 & 0.619 & 0.718 & 0.714 \\
        5  & 0.681 & 0.671 & 0.736 & 0.716 \\
        10 & 0.673 & 0.667 & 0.739 & 0.718 \\
        20 &  0.680 & 0.670 & 0.746 & 0.731 \\
        \hline
    \end{tabular}
\end{table}

\begin{table}[htbp]
    \centering
    \caption{Experiment D (IEMOCAP): SNR ablation with natural soundscapes for ResNet-18 and AST. Metrics reported are accuracy (Acc.) and macro-F1 per SNR.}
    \label{tab:snr_ablation_iemocap}
    \begin{tabular}{lcccc}
        \hline
        SNR (dB) & ResNet-18 Acc. & ResNet-18 Macro-F1 & AST Acc. & AST Macro-F1 \\
        \hline
        0  & 0.570 & 0.406 & 0.615 & 0.448 \\
        5  & 0.577 & 0.366 & 0.624 & 0.460 \\
        10 & 0.596 & 0.370 & 0.642 & 0.526 \\
        20 & 0.609 & 0.351 & 0.636 & 0.457 \\
        \hline
    \end{tabular}
\end{table}








\begin{figure}[htbp]
    \centering
    \begin{tikzpicture}
    \pgfplotsset{compat=1.18}

    \begin{groupplot}[
      group style={group size=2 by 1, horizontal sep=3.5em},
      width=0.46\textwidth, height=0.36\textwidth,
      xmin=-0.5, xmax=20.5,
      ymin=0.30, ymax=0.80,
      xtick={0,5,10,20},
      xlabel={SNR (dB)}, ylabel={Macro-F1},
      grid=both,
      grid style={line width=.1pt, draw=gray!30},
      major grid style={line width=.2pt, draw=gray!55},
      ticklabel style={font=\small},
      label style={font=\small},
      ytick distance=0.1
    ]

    % ---------- RAVDESS ----------
    \nextgroupplot[title={RAVDESS}, title style={font=\small},
      legend style={nodes={scale=0.9, transform shape}, at={(0.5,1.25)}, anchor=south, legend columns=2}
    ]
    \addplot+[very thick, color=orange!95!black, mark=*, mark size=2.5pt, mark options={fill=orange!80!white}]
      coordinates {(0,0.619) (5,0.671) (10,0.667) (20,0.670)};
    \addlegendentry{ResNet-18}

    \addplot+[very thick, color=violet!90!black, mark=triangle*, mark size=3pt, mark options={fill=violet!80!white}]
      coordinates {(0,0.714) (5,0.716) (10,0.718) (20,0.731)};
    \addlegendentry{AST Base}

    % ---------- IEMOCAP ----------
    \nextgroupplot[title={IEMOCAP}, title style={font=\small}]
    \addplot+[very thick, color=orange!95!black, mark=*, mark size=2.5pt, mark options={fill=orange!80!white}]
      coordinates {(0,0.406) (5,0.366) (10,0.370) (20,0.351)};
    \addplot+[very thick, color=violet!90!black, mark=triangle*, mark size=3pt, mark options={fill=violet!80!white}]
      coordinates {(0,0.448) (5,0.460) (10,0.526) (20,0.457)};

    \end{groupplot}
    \end{tikzpicture}
    \caption{Experiment D: Macro-F1 vs. SNR with natural soundscapes for ResNet-18 and AST Base.}
    \label{fig:snr_ablation_macroF1}
\end{figure}